\section{Higher direct images of coherent sheaves}
\label{section-proper-pushforward}

\noindent
In this section we prove the fundamental fact that the higher
direct images of a coherent sheaf under a proper morphism
are coherent.

\begin{proposition}
\label{proposition-proper-pushforward-coherent}
\begin{reference}
\cite[III Theorem 3.2.1]{EGA}
\end{reference}
Let $S$ be a locally Noetherian scheme.
Let $f : X \to S$ be a proper morphism.
Let $\mathcal{F}$ be a coherent $\mathcal{O}_X$-module.
Then $R^if_*\mathcal{F}$ is a coherent $\mathcal{O}_S$-module
for all $i \geq 0$.
\end{proposition}

\begin{proof}
Since the problem is local on $S$ we may assume that $S$ is
a Noetherian scheme. Since a proper morphism is of finite type
we see that in this case $X$ is a Noetherian scheme also.
Consider the property $\mathcal{P}$ of coherent sheaves
on $X$ defined by the rule
$$
\mathcal{P}(\mathcal{F}) \Leftrightarrow
R^pf_*\mathcal{F}\text{ is coherent for all }p \geq 0
$$
We are going to use the result of
Lemma \ref{lemma-property} to prove that
$\mathcal{P}$ holds for every coherent sheaf on $X$.

\medskip\noindent
Let
$$
0 \to \mathcal{F}_1 \to \mathcal{F}_2 \to \mathcal{F}_3 \to 0
$$
be a short exact sequence of coherent sheaves on $X$.
Consider the long exact sequence of higher direct images
$$
R^{p - 1}f_*\mathcal{F}_3 \to
R^pf_*\mathcal{F}_1 \to
R^pf_*\mathcal{F}_2 \to
R^pf_*\mathcal{F}_3 \to
R^{p + 1}f_*\mathcal{F}_1
$$
Then it is clear that if 2-out-of-3 of the sheaves $\mathcal{F}_i$
have property $\mathcal{P}$, then the higher direct images of the
third are sandwiched in this exact complex between two coherent
sheaves. Hence these higher direct images are also coherent by
Lemma \ref{lemma-coherent-abelian-Noetherian} and
\ref{lemma-coherent-Noetherian-quasi-coherent-sub-quotient}.
Hence property $\mathcal{P}$ holds for the third as well.

\medskip\noindent
Let $Z \subset X$ be an integral closed subscheme.
We have to find a coherent sheaf $\mathcal{F}$ on $X$ whose support is
contained in $Z$, whose stalk at the generic point $\xi$ of $Z$ is a
$1$-dimensional vector space over $\kappa(\xi)$ such that $\mathcal{P}$
holds for $\mathcal{F}$. Denote $g = f|_Z : Z \to S$ the restriction of $f$.
Suppose we can find a coherent sheaf $\mathcal{G}$ on $Z$ such
that
(a) $\mathcal{G}_\xi$ is a $1$-dimensional vector space over $\kappa(\xi)$,
(b) $R^pg_*\mathcal{G} = 0$ for $p > 0$, and
(c) $g_*\mathcal{G}$ is coherent. Then we can consider
$\mathcal{F} = (Z \to X)_*\mathcal{G}$. As $Z \to X$ is a closed immersion
we see that $(Z \to X)_*\mathcal{G}$ is coherent on $X$
and $R^p(Z \to X)_*\mathcal{G} = 0$ for $p > 0$
(Lemma \ref{lemma-finite-pushforward-coherent}).
Hence by the relative Leray spectral sequence
(Cohomology, Lemma \ref{cohomology-lemma-relative-Leray})
we will have $R^pf_*\mathcal{F} = R^pg_*\mathcal{G} = 0$ for $p > 0$
and $f_*\mathcal{F} = g_*\mathcal{G}$ is coherent.
Finally $\mathcal{F}_\xi = ((Z \to X)_*\mathcal{G})_\xi = \mathcal{G}_\xi$
which verifies the condition on the stalk at $\xi$.
Hence everything depends on finding a coherent sheaf $\mathcal{G}$
on $Z$ which has properties (a), (b), and (c).

\medskip\noindent
We can apply Chow's Lemma \ref{lemma-chow-Noetherian}
to the morphism $Z \to S$. Thus we get a diagram
$$
\xymatrix{
Z \ar[rd]_g & Z' \ar[d]^-{g'} \ar[l]^\pi \ar[r]_i & \mathbf{P}^m_S \ar[dl] \\
& S &
}
$$
as in the statement of Chow's lemma. Also, let $U \subset Z$ be
the dense open subscheme such that $\pi^{-1}(U) \to U$ is an isomorphism.
By the discussion in Remark \ref{remark-chow-Noetherian} we see that
$i' = (i, \pi) : Z' \to \mathbf{P}^m_Z$ is
a closed immersion. Hence
$$
\mathcal{L} = i^*\mathcal{O}_{\mathbf{P}^m_S}(1) \cong
(i')^*\mathcal{O}_{\mathbf{P}^m_Z}(1)
$$
is $g'$-relatively ample and $\pi$-relatively ample (for example by
Morphisms, Lemma \ref{morphisms-lemma-characterize-ample-on-finite-type}).
Hence by Lemma \ref{lemma-kill-by-twisting}
there exists an $n \geq 0$ such that
both $R^p\pi_*\mathcal{L}^{\otimes n} = 0$ for all $p > 0$ and
$R^p(g')_*\mathcal{L}^{\otimes n} = 0$ for all $p > 0$.
Set $\mathcal{G} = \pi_*\mathcal{L}^{\otimes n}$.
Property (a) holds because $\pi_*\mathcal{L}^{\otimes n}|_U$ is
an invertible sheaf (as $\pi^{-1}(U) \to U$ is an isomorphism).
Properties (b) and (c) hold because by the relative Leray
spectral sequence
(Cohomology, Lemma \ref{cohomology-lemma-relative-Leray})
we have
$$
E_2^{p, q} = R^pg_* R^q\pi_*\mathcal{L}^{\otimes n}
\Rightarrow
R^{p + q}(g')_*\mathcal{L}^{\otimes n}
$$
and by choice of $n$ the only nonzero terms in $E_2^{p, q}$ are
those with $q = 0$ and the only nonzero terms of
$R^{p + q}(g')_*\mathcal{L}^{\otimes n}$ are those with $p = q = 0$.
This implies that $R^pg_*\mathcal{G} = 0$ for $p > 0$ and
that $g_*\mathcal{G} = (g')_*\mathcal{L}^{\otimes n}$.
Finally, applying the previous
Lemma \ref{lemma-locally-projective-pushforward}
we see that $g_*\mathcal{G} = (g')_*\mathcal{L}^{\otimes n}$ is
coherent as desired.
\end{proof}

\begin{lemma}
\label{lemma-proper-over-affine-cohomology-finite}
Let $S = \Spec(A)$ with $A$ a Noetherian ring.
Let $f : X \to S$ be a proper morphism.
Let $\mathcal{F}$ be a coherent $\mathcal{O}_X$-module.
Then $H^i(X, \mathcal{F})$ is a finite $A$-module for all $i \geq 0$.
\end{lemma}

\begin{proof}
This is just the affine case of
Proposition \ref{proposition-proper-pushforward-coherent}.
Namely, by Lemmas \ref{lemma-quasi-coherence-higher-direct-images} and
\ref{lemma-quasi-coherence-higher-direct-images-application} we know that
$R^if_*\mathcal{F}$ is the quasi-coherent sheaf associated
to the $A$-module $H^i(X, \mathcal{F})$
and by Lemma \ref{lemma-coherent-Noetherian} this is
a coherent sheaf if and only if $H^i(X, \mathcal{F})$
is an $A$-module of finite type.
\end{proof}

\begin{lemma}
\label{lemma-graded-finiteness}
Let $A$ be a Noetherian ring.
Let $B$ be a finitely generated graded $A$-algebra.
Let $f : X \to \Spec(A)$ be a proper morphism.
Set $\mathcal{B} = f^*\widetilde B$.
Let $\mathcal{F}$ be a quasi-coherent
graded $\mathcal{B}$-module of finite type.
\begin{enumerate}
\item For every $p \geq 0$ the graded $B$-module $H^p(X, \mathcal{F})$
is a finite $B$-module.
\item If $\mathcal{L}$ is an ample invertible $\mathcal{O}_X$-module,
then there exists an integer $d_0$ such that
$H^p(X, \mathcal{F} \otimes \mathcal{L}^{\otimes d}) = 0$
for all $p > 0$ and $d \geq d_0$.
\end{enumerate}
\end{lemma}

\begin{proof}
To prove this we consider the fibre product diagram
$$
\xymatrix{
X' = \Spec(B) \times_{\Spec(A)} X
\ar[r]_-\pi \ar[d]_{f'} &
X \ar[d]^f \\
\Spec(B) \ar[r] &
\Spec(A)
}
$$
Note that $f'$ is a proper morphism, see
Morphisms, Lemma \ref{morphisms-lemma-base-change-proper}.
Also, $B$ is a finitely generated $A$-algebra, and hence
Noetherian (Algebra, Lemma \ref{algebra-lemma-Noetherian-permanence}).
This implies that $X'$ is a Noetherian scheme
(Morphisms, Lemma \ref{morphisms-lemma-finite-type-noetherian}).
Note that $X'$ is the relative spectrum of the quasi-coherent
$\mathcal{O}_X$-algebra $\mathcal{B}$ by
Constructions, Lemma \ref{constructions-lemma-spec-properties}.
Since $\mathcal{F}$ is a quasi-coherent $\mathcal{B}$-module
we see that there is a unique quasi-coherent
$\mathcal{O}_{X'}$-module $\mathcal{F}'$ such that
$\pi_*\mathcal{F}' = \mathcal{F}$, see
Morphisms, Lemma \ref{morphisms-lemma-affine-equivalence-modules}
Since $\mathcal{F}$ is finite type as a $\mathcal{B}$-module we
conclude that $\mathcal{F}'$ is a finite type
$\mathcal{O}_{X'}$-module (details omitted). In other words,
$\mathcal{F}'$ is a coherent $\mathcal{O}_{X'}$-module
(Lemma \ref{lemma-coherent-Noetherian}).
Since the morphism $\pi : X' \to X$ is affine we have
$$
H^p(X, \mathcal{F}) = H^p(X', \mathcal{F}')
$$
by Lemma \ref{lemma-relative-affine-cohomology}.
Thus (1) follows from
Lemma \ref{lemma-proper-over-affine-cohomology-finite}.
Given $\mathcal{L}$ as in (2) we set
$\mathcal{L}' = \pi^*\mathcal{L}$. Note that $\mathcal{L}'$ is
ample on $X'$ by
Morphisms, Lemma \ref{morphisms-lemma-pullback-ample-tensor-relatively-ample}.
By the projection formula
(Cohomology, Lemma \ref{cohomology-lemma-projection-formula}) we have
$\pi_*(\mathcal{F}' \otimes \mathcal{L}') = \mathcal{F} \otimes \mathcal{L}$.
Thus part (2) follows by the same reasoning as above from
Lemma \ref{lemma-kill-by-twisting}.
\end{proof}












